\documentclass[a4paper]{article}

\usepackage[default]{lato}
\usepackage[swedish]{babel}
\usepackage[dvipsnames]{xcolor}
\usepackage[T1]{fontenc}
\usepackage[utf8]{inputenc}
\usepackage{graphicx}
\usepackage{fancyhdr}
\usepackage{lastpage}
\usepackage{parskip}

\definecolor{cerise}{RGB}{232, 61, 132}
\definecolor{gray3}{RGB}{66, 66, 66}

\addtolength{\textheight}{4cm}
\addtolength{\voffset}{-2cm}
\addtolength{\textwidth}{2cm}
\addtolength{\hoffset}{-1cm}
\setlength{\headheight}{40pt}

\pagestyle{fancy}
\thispagestyle{fancy}
\fancyhead[L]{\small{
    Motion till Val-SM % Skriv SM-namn
\\\textcolor{gray3}{
    15-05-2023 % Datum för SM
}}}
\fancyhead[C]{\includegraphics[width=1.25cm,height=1.25cm]{./skold-color.pdf}}
\fancyhead[R]{\small{\textcolor{gray3}{
    Erik Hedlund\\ % Skriv ditt namn här
    sida \thepage{} av \pageref{LastPage}}}}
\fancyfoot[C]{\small{\textcolor{gray3}{Konglig Datasektionen, THS 100 44, datasektionen.se}}}
\renewcommand{\headrulewidth}{0pt}

\title{\textcolor{cerise}{\textbf{
    Motion angående\\
    Återkommande projekt i reglementet
}}}

\date{}

\begin{document}

\maketitle
\thispagestyle{fancy} % Smutsigt fulhack för att få header på första sidan

\section*{\textcolor{cerise}{
    Bakgrund
}}

Sektionen har idag ett flertal återkommande projekt som engagerar många av dess medlemmar.
Att dessa projekt regelbundna aktivitet dock endast sätts av praxis och inte finns inte
representerat i våra styrdokument är inte optimalt då det dels gör att man inte kan
hitta information som när projekt öppnas, när projektledare väljs samt vad projektledare benämns.
Annan problematik det här medför är också att det blir omöjligt att syfta på projekten samt deras
projektledare under andra paragrafer i reglementet, ett exempel på detta är att det är omöjligt
att göra projektledarroller till intervjuposter. Ett sista problem med projektens frånvarande
från reglementet är att D-rektoratet kan göra frikostiga förändringar så som att byta namn
samt på projekt, flytta verksamhetsperiod eller till med avveckla utan att SM har möjlighet att påverka.

Lösningen på återkommande projekt-``problematiken'' har traditionellt varit att göra om dessa
projekt till nämnder. Denna lösning är inte optimal för många av projekten vi har som återkommande
nu, då detta innebär en mängd oavsiktliga förändringar som påverkar projekten negativt, så som
att budget sätts efter kalenderår istället för projektperiod samt att det rent formellt gör det
omöjligt för två iteration av ett projekt att löpa parallellt.

Förändringarna i den här motionen skulle inte innebära några större förändringar i praktiken, då
mycket redan efterlevs genom praxis, som ovan nämnt. Den stora praktiska förändringen med den här
motionen kommer vara att återkommande projekt öppnas utan att det måste beslutas på DM.
Det överhängande syftet med motionen är att möjliggöra att i framtiden syfta på dessa projekt i
andra reglementesparagrafer.

\section*{\textcolor{cerise}{
    Förslag till beslut
}}

	Mot bakgrund av ovanstående yrkar jag:

\newcounter{attnummer}
\setcounter{attnummer}{1}
\begin{list}{\bf att$_{\theattnummer}$\stepcounter{attnummer}}{}

% Stadgeändring

\item I stadgarna ändra \textit{\S7.2 Uppstart av återkommande projekt} från
    \begin{quote}
        Projekt som D-rektoratet anser är regelbundet återkommande kan startas utan SM-beslut. De startas då istället genom beslut på DM med en kort motivering innehållande referens till minst ett väldigt likartat tidigare projekt samt en uppskattning av de uppgifter, bortsett från budget, som krävs för att starta ett nytt projekt. Val av projektledare ska ske på nästkommande SM enligt samma procedur som för övriga funktionärsposter. Vald projektledare åläggs att inkomma med motion innehållandes budget samt verksamhetsplan för projektet till första möjliga SM efter valet, såvida SM inte redan beslutat om dessa för denna projektomgång. D-rektoratet ansvarar för att en förteckning över återkommande projekt finns tillgänglig för alla sektionens medlemmar på dess hemsida.
    \end{quote}
    till
    \begin{quote}
        Projekt skrivna som återkommande under reglementet startas utan beslut på bestämt datum. Val av projektledare ska ske på nästkommande SM enligt samma procedur som för övriga funktionärsposter. Om D-rektoratet anser att ett återkommande projekt inte ska öppnas för nästa verksamhetsperiod så kan de bestämma det på ett DM innan datumet som projektet öppnas.
    \end{quote}

% Basic saker

\item I reglementet tillägga \textit{\S4 Återkommande projekt} med text
    \begin{quote}
        \textit{Återkommande projekt} är projekt som öppnas årligen under under bestämt datum.
        Projektledare för återkommande projekt väljs på samma sätt som övriga projektledare
        nästkommande SM efter projektets öppnande.

        Vald projektledare åläggs att inkomma med motion innehållande budget samt
        verksamhetsplan för projektet till första möjliga SM efter valet, såvida SM inte
        redan beslutat om dessa för denna projektomgång.

        Återkommande projekt har slutgiltiga mål där projekten planeras bedriva verksamhet framtill
        att dessa är uppfyllda. Att fortsätta verksamhet efter dessa mål uppfyllts,
        alternativt upphöra med verksamhet innan dessa mål uppfyllts, är möjligt om skäl för detta
        finns.
    \end{quote}

% METAspexet

\item I reglementet tillägga \textit{\S4.1 METAspexet}
\item I reglementet tillägga \textit{\S4.1.1 Ändamål} med text
    \begin{quote}
        METAspexet är ett återkommande projekt som har delad verksamhet med motsvarande nämnd på Sektionen för Medieteknik

        METAspexet skall verka för att ge sektionsmedlemmar av Konglig Datasektionen och Sektionen för Medieteknik chans att utveckla flertalet färdigheter relaterad till teaterproduktion samt för att förstärka relationen mellan Konglig Datasektionen och Sektionen för Medieteknik.
    \end{quote}
\item I reglementet tillägga \textit{\S4.1.2 Organisation} med text
    \begin{quote}
        METAspexet leds av projektledare, benämnd Direqteur, för Konglig Datasektionen samt motsvarande nämndordförande
        för Sektionen för Medieteknik. Dessa projektledare och nämndordförande benämns \textit{METAspexets direqteurer}

        METAspexets direqteurer ansvarar för att rekrytera samt leda \textit{METAspexets chefsgrupp}

        METAspexet chefsgrupp ansvarar för att rekrytera samt leda METAspexet resterande grupper.
    \end{quote}
\item I reglementet tillägga \textit{\S4.1.3 Verksamhet} med text
    \begin{quote}
        \begin{itemize}
            \item Att sätta upp spexföreställningar under nästkommande kalenderår.
            \item Att ge medlemmar utrymme till engagemang och kreativt arbete samtidigt som de känner att de bidrar till ett större projekt.
            \item Att ge medlemmar möjlighet att lära sig nya saker tillsammans och av varandra.
            \item Att värna och utveckla relationen mellan Data och Media som sektioner, såväl som relationerna mellan sektionernas medlemmar.
        \end{itemize}
    \end{quote}
\item I reglementet tillägga \textit{\S4.1.4 Period} med text
    \begin{quote}
        Nya METAspexet-projekt öppnas årligen 1:a april och bör bedriva verksamhet framtill att dess alla föreställningar har skett nästkommande kalenderår.

        Projekttitel för respektive årgång är \textit{METAspexet <år>}, där <år> är det år METAspexets föreställningar planeras hållas.
    \end{quote}

% dÅre

\item I reglementet tillägga \textit{\S4.2 dåre}
\item I reglementet tillägga \textit{\S4.2.1 Ändamål} med text
    \begin{quote}
        dÅre är ett återkommande projekt som syftar att främja idrottsengagemanget inom sektionen samt främja gemenskap mellan studenter från olika årgångar genom att organisera en skidsemester.
    \end{quote}
\item I reglementet tillägga \textit{\S4.2.2 Organisation} med text
    \begin{quote}
        dÅre leds av projektledare

        Projektledare för dÅre ansvarar för att rekrytera samt leda dÅres projektgrupp.
    \end{quote}
\item I reglementet tillägga \textit{\S4.2.3 Verksamhet} med text
    \begin{quote}
        \begin{itemize}
            \item Att organisera en skidsemester för Konglig Datasektionens medlemmar.
            \item Att uppmuntra medlemmar oberoende av tidigare erfarenhet att delta i alpint idrottsengagemanget.
        \end{itemize}
    \end{quote}
\item I reglementet tillägga \textit{\S4.2.4 Period} med text
    \begin{quote}
        Nya dåre-projekt öppnas årligen 1:a april och bör bedriva verksamhet framtill att dess planerade skidresa skett nästkommande kalenderår.

        Projekttitel för respektive årgång är \textit{dÅre <år>}, där <år> är det år dÅres skidresa planeras hållas.
    \end{quote}

% Studs

\item I reglementet tillägga \textit{\S4.3 Studs}
\item I reglementet tillägga \textit{\S4.3.1 Ändamål} med text
    \begin{quote}
        TODO
    \end{quote}
\item I reglementet tillägga \textit{\S4.3.2 Organisation} med text
    \begin{quote}
        Studs leds av projektledare

        Två projektledare bör väljas

        Projektledare för Studs ansvarar för att rekrytera samt leda Studs \textit{ledningsgrupp}.

        Ledningsgruppen ansvarar för att rekrytera samt leda Studs övriga grupper.
    \end{quote}
\item I reglementet tillägga \textit{\S4.3.3 Verksamhet} med text
    \begin{quote}
        TODO
    \end{quote}
\item I reglementet tillägga \textit{\S4.3.4 Period} med text
    \begin{quote}
        Nya Studs-projekt öppnas årligen 1:a april och bör bedriva verksamhet framtill att dess studieresa har skett nästkommande kalenderår.

        Projekttitel för respektive årgång är \textit{Studs <år>}, där <år> är det år Studs resa planeras hållas.
    \end{quote}

% Vårbalen

\item I reglementet tillägga \textit{\S4.4 Vårbalen}
\item I reglementet tillägga \textit{\S4.4.1 Ändamål} med text
    \begin{quote}
        Vårbalen är ett återkommande projekt med syfte att anordna en finsittning under våren.
    \end{quote}
\item I reglementet tillägga \textit{\S4.4.2 Organisation} med text
    \begin{quote}
        Vårbalen leds av projektledare, benämnd \textit{Vårbalsgeneral}.

        Vårbalsgeneral ansvarar för att rekrytera samt leda Vårbalens projektgrupp.
    \end{quote}
\item I reglementet tillägga \textit{\S4.4.3 Verksamhet} med text
    \begin{quote}
        \begin{itemize}
            \item Anordna en finsittning med tillhörande efterkör under våren.
            \item Bjuda in externa representanter från exempelvis andra sektioner inom THS till Vårbalen.
            \item Bjuda in Ordförande Emeriti till Vårbalen.
        \end{itemize}
    \end{quote}
\item I reglementet tillägga \textit{\S4.4.4 Period} med text
    \begin{quote}
        Nya Vårbalen-projekt öppnas årligen 1:a april och bör bedriva verksamhet framtill att dess planerade sittning skett nästkommande kalenderår.

        Projekttitel för respektive årgång är \textit{Vårbalen <år>}, där <år> är det år Vårbalens sittning planeras hållas.
    \end{quote}

% Project pride

\item I reglementet tillägga \textit{\S4.5 Project pride}
\item I reglementet tillägga \textit{\S4.5.1 Ändamål} med text
    \begin{quote}
        Project Pride är ett återkommande projekt som ämnar organisera
        deltagande i Stockholm Pride Parade, samt planera event under våren.
        Project Pride verkar för att uppmärksamma och främja HBTQIA+.
    \end{quote}
\item I reglementet tillägga \textit{\S4.5.2 Organisation} med text
    \begin{quote}
        Project Pride leds av projektledare, benämnd \textit{Ansvarig för Project Pride}.

        Projektledare för Project Pride ansvarar för att rekrytera samt leda projektgruppen
        för Project Pride.
    \end{quote}
\item I reglementet tillägga \textit{\S4.5.3 Verksamhet} med text
    \begin{quote}
        \begin{itemize}
            \item Arbeta för att uppmärksamma och främja Konglig Datasektionens JML-arbete, främst mot HBTQIA+, både internt och externt.
            \item Anordna ett bidrag till Stockholm Pride Parade 2023 för Konglig Data-sektionen, samt arbeta för att anordna fler evenemang under våren 2023.
            \item Möjliggöra att så många av Konglig Datasektionens medlemmar som möjligt ska kunna delta under evenemangen.
            \item Se om möjlighet finns till samarbete med andra delar av sektionen, och om så finnes arbeta för att främja detta.
        \end{itemize}
    \end{quote}
\item I reglementet tillägga \textit{\S4.5.4 Period} med text
    \begin{quote}
        Nya Project Pride-projekt öppnas årligen 16:e november och bör bedriva verksamhet framtill att dess ??? har skett nästkommande kalenderår.

        Projekttitel för respektive årgång är Project Pride <år>, där <år> är det år Project Prides deltagande i Stockholm Pride Parade planeras hållas.
    \end{quote}

% Funktionärer

\item I reglementet tillägga \textit{\S5.3 Projektledare för återkommande projekt}
\item I reglementet tillägga \textit{\S5.3.1 Direqteur} med text
    \begin{quote}
        Är projektledare för METAspexet
    \end{quote}
\item I reglementet tillägga \textit{\S5.3.2 Projektledare för dÅre} med text
    \begin{quote}
        Är projektledare för dÅre
    \end{quote}
\item I reglementet tillägga \textit{\S5.3.4 Vårbalsgeneral} med text
    \begin{quote}
        Är projektledare för Vårbalen
    \end{quote}
\item I reglementet tillägga \textit{\S5.3.5 Ansvarig för Project Pride} med text
    \begin{quote}
        Är projektledare för Project Pride
    \end{quote}

% Fippel

\item Konsekvensnumrera alla punkter jag fuckar upp genom att peta in Återkommande projekt under \S4.
\end{list}

\vspace{2cm}
\noindent
% Name, eventuell titel eller annat skoj
Erik Hedlund, Projektledare för ett återkommande projekt

% Avkommentera följande om du vill visa en bild
\begin{center}
\includegraphics[scale=0.25]{
    % Eventuell fin bild du vill visa.
    ./meme.png
}
\end{center}


\end{document}
